% ==============================================================================
% INTRODUCAO — reestruturada em 30/10/2026 conforme o roteiro do orientador:
% motivacao do problema; importancia dos modelos preditivos para a MSC; o que
% a literatura tem feito; desafios e lacunas; objetivo geral e especificos.
% A escolha dos dados e a delimitacao da coorte foram para a metodologia
% (secao Populacao de estudo). Versao anterior: historico do git (versao 27).
% ==============================================================================

\chapter{Introdução}
\label{cap:introducao}

\section{Motivação do problema}
\label{sec:intro-motivacao}

A insuficiência cardíaca é uma síndrome crônica e progressiva em que o
coração não consegue manter o débito adequado às necessidades do organismo.
Entre os pacientes acometidos, a morte ocorre por duas vias principais: a
falência de bomba, em que a função cardíaca se deteriora progressivamente, e
a \gls{MSC}, um evento abrupto, em geral causado por uma arritmia ventricular
--- taquicardia ou fibrilação ventricular --- em um paciente que, muitas
vezes, se encontrava estável. A \gls{MSC} responde por parcela expressiva dos
óbitos nessa população, estimada em cerca de metade das mortes de pacientes
com insuficiência cardíaca \cite{tomaselli2004sudden}.

A gravidade do problema decorre também do caráter inesperado do evento.
Diferentemente da falência de bomba, que costuma ser precedida de piora
clínica e de internações, a \gls{MSC} pode ocorrer sem sinais de alerta, e a
chance de sobrevivência depende de atendimento em poucos minutos. A
prevenção, portanto, precisa ser feita antes do evento, identificando os
pacientes em risco. A principal medida disponível é o
cardiodesfibrilador implantável, que reverte a arritmia no momento em que
ela ocorre e reduziu a mortalidade de pacientes com insuficiência cardíaca e
fração de ejeção reduzida \cite{bardy2005amiodarone}.

A indicação do dispositivo para prevenção primária baseia-se,
principalmente, na fração de ejeção do ventrículo esquerdo
\cite{zeppenfeld2022esc}. Esse critério, no entanto, é reconhecidamente
imperfeito: parte dos pacientes que recebem o dispositivo nunca chega a
utilizá-lo, enquanto parte das mortes súbitas ocorre em pacientes com fração
de ejeção acima do limiar de indicação \cite{zeppenfeld2022esc}. Há,
portanto, espaço para ferramentas que estratifiquem melhor o risco de
\gls{MSC}, distinguindo-o do risco de morte por falência de bomba.

\section{Modelos preditivos para a morte súbita cardíaca}
\label{sec:intro-modelos}

Um modelo preditivo de \gls{MSC} estima, para cada paciente, a probabilidade
de que ele venha a sofrer o evento em um horizonte de tempo definido. Esse
tipo de modelo permite ordenar os pacientes por risco e concentrar a
prevenção naqueles com maior probabilidade de benefício, em vez de
depender de um único limiar clínico. Para que seja útil, o modelo precisa
basear-se em informações disponíveis na rotina e ser avaliado de forma que
seu desempenho em pacientes novos possa ser estimado com confiança.

O \gls{ECG} de Holter de 24~horas reúne essas condições. É um exame não
invasivo, de baixo custo e amplamente disponível, que registra a atividade
elétrica do coração ao longo de um dia inteiro, captando batimentos
ectópicos, pausas e variações da frequência cardíaca entre o sono e a
vigília que um traçado de repouso de poucos segundos não alcança. Como a
\gls{MSC} tem origem, na maior parte dos casos, em uma instabilidade elétrica,
é plausível que o Holter contenha sinais do substrato arrítmico que a
antecede. O aprendizado de máquina oferece meios para explorar esses sinais,
seja a partir de medidas extraídas do traçado, seja diretamente a partir do
sinal bruto.

\section{O que a literatura tem feito}
\label{sec:intro-literatura}

Os estudos de predição de \gls{MSC} a partir do \gls{ECG} ambulatorial têm
seguido dois caminhos, descritos em detalhe no
\autoref{cap:trabalhos-relacionados}. O primeiro, mais antigo, extrai do
Holter marcadores definidos por especialistas --- variabilidade da frequência
cardíaca, ectopia, turbulência da frequência, alternância da onda T,
dinâmica da repolarização --- e os combina com variáveis clínicas em modelos
estatísticos ou de aprendizado de máquina clássico. Na base \gls{MUSIC},
utilizada neste trabalho, esse caminho produziu escores de risco que
combinam marcadores da repolarização e variáveis clínicas
\cite{ramirez2017sudden} e modelos de aprendizado de máquina sobre
marcadores de segmentos curtos do Holter \cite{swee2026multimodal}.

O segundo caminho, mais recente, aplica redes neurais profundas diretamente
ao sinal, apoiado no sucesso dessas redes em tarefas diagnósticas sobre o
\gls{ECG}, como a detecção de arritmias no Holter
\cite{hannun2019cardiologist} e o diagnóstico automático no \gls{ECG} de 12
derivações \cite{ribeiro2020automatic}. Na predição de eventos cardíacos a
partir do Holter, as abordagens que obtiveram melhores resultados partiram
de redes pré-treinadas em grandes bases de \gls{ECG} e ajustaram apenas um
classificador pequeno com os poucos eventos disponíveis \cite{kim2025deep}.

\section{Desafios e lacunas}
\label{sec:intro-lacunas}

A predição de \gls{MSC} difere das tarefas diagnósticas em que as redes
profundas se destacaram. Nas tarefas diagnósticas, o padrão procurado está
presente no traçado no momento do exame e há dezenas de milhares de
exemplos rotulados. Na predição de \gls{MSC}, o evento ocorre meses ou anos
depois do Holter, e não se sabe se o registro daquele dia já contém algum
sinal do que vai acontecer; além disso, as coortes com causa de óbito
adjudicada reúnem poucas centenas de pacientes. Nessas condições, a escolha
entre aprender diretamente do sinal ou de medidas extraídas dele não é óbvia
e precisa ser avaliada empiricamente.

Três lacunas motivam este trabalho. A primeira é que os estudos anteriores
avaliaram, cada um, uma única abordagem, em coortes e com protocolos
diferentes, o que impede saber qual delas prevê melhor a \gls{MSC}. Faltam
comparações diretas, nos mesmos pacientes e com a mesma forma de avaliação,
entre redes profundas aplicadas ao sinal bruto, modelos baseados em
marcadores extraídos do \gls{ECG} e modelos baseados em variáveis clínicas.

A segunda é que parte dos estudos adota escolhas de avaliação que inflam o
desempenho estimado. Uma delas é calcular as métricas sobre trechos do
sinal, e não sobre pacientes: como um mesmo exame fornece centenas de
trechos, cada paciente passa a contar centenas de vezes. Outra é selecionar
as variáveis observando o desfecho de toda a coorte, inclusive dos pacientes
usados no teste \cite{kapoor2023leakage}. Neste trabalho, os trechos são
usados no treinamento, mas toda avaliação é feita por paciente.

A terceira é a inclusão de pacientes com fibrilação atrial, ritmo
facilmente reconhecível no sinal e associado a pior prognóstico, que pode
elevar a discriminação sem que o modelo capte marcadores de vulnerabilidade
elétrica ventricular \cite{swee2026multimodal}.

\section{Objetivos}
\label{sec:intro-objetivo}

\subsection{Objetivo geral}
\label{subsec:intro-objetivo-geral}

Este trabalho propõe a investigação de métodos baseados em aprendizado de
máquina e no pré-processamento do sinal de \gls{ECG} para a predição da
ocorrência de \gls{MSC} em uma população de pacientes com insuficiência
cardíaca crônica, a partir do Holter de 24~horas, ao longo de um período de
acompanhamento de quatro anos.

\subsection{Objetivos específicos}
\label{subsec:intro-objetivos-especificos}

\begin{enumerate}
	\item[(i)] Comparar três cenários para a predição da \gls{MSC}: (1) modelos de
	      aprendizado profundo aplicados ao sinal bruto do \gls{ECG}; (2)
	      pré-processamento do sinal, extração de características e uso de
	      modelos de aprendizado de máquina clássico; e (3) modelos baseados
	      em variáveis clínicas, isoladas e combinadas aos dados extraídos do
	      \gls{ECG}.
	\item[(ii)] Avaliar a combinação de modelos baseados em dados clínicos com
	      modelos baseados em dados extraídos do \gls{ECG}, quantificando o
	      ganho de desempenho de cada combinação.
	\item[(iii)] Revisar os trabalhos da literatura que propõem a predição de
	      \gls{MSC} a partir do \gls{ECG} ambulatorial e posicionar os
	      resultados deste trabalho em relação a eles.
\end{enumerate}

\section{Organização do trabalho}
\label{sec:intro-organizacao}

O \autoref{cap:trabalhos-relacionados} apresenta os trabalhos relacionados.
O capítulo de fundamentação teórica descreve o pré-processamento e a
extração de características do \gls{ECG}, os modelos de aprendizado profundo
e de aprendizado de máquina aplicados ao sinal e os métodos de avaliação. A
metodologia apresenta a população de estudo e detalha os três cenários, na
mesma ordem em que seus resultados são apresentados no capítulo seguinte. O
último capítulo reúne as conclusões, as limitações e os trabalhos futuros.
